\documentclass[10pt]{article}
\usepackage[margin=0.7in]{geometry}
\usepackage{booktabs}
\usepackage{graphicx}

\title{COS568 Learned Index\\Milestone 2 Report}
\author{Liv d'Aliberti}
\date{\today}

\begin{document}
\maketitle
\thispagestyle{empty}

\section*{Setup}

\paragraph{Workflow.}
I set up a dedicated workflow that mirrors Part~1: dataset check/download, workload generation, benchmark build, mixed-workload benchmarking, plot generation, and PDF report assembly. Each benchmark configuration is repeated three times (\texttt{-r 3}), and the reporting scripts choose the row with the highest mean mixed-workload throughput for each index and workload. The Slurm pipeline remains the default reproducible path, and the Facebook report artifacts below were generated with the same scripts and repository \texttt{.venv} while the queued cluster job was still pending on resources.

\paragraph{Hybrid Designs.}
I implemented two hybrid variants. The required baseline is \texttt{HybridPGMLIPP}, a naive synchronous design: the bulk-loaded base keys go into LIPP, incoming inserts go into a Dynamic PGM buffer, lookups probe the buffer first and then fall back to LIPP, and the buffered keys are flushed into LIPP once the buffer reaches a configured threshold. I also implemented \texttt{HybridPGMLIPPIncremental}, a less naive incremental-draining design that rotates a full Dynamic PGM buffer into a draining state, keeps a fresh Dynamic PGM buffer available for new inserts, and migrates a bounded number of keys from the draining buffer into LIPP on subsequent insert operations.

\paragraph{Hyperparameter Sweep.}
For Facebook, I sweep Dynamic PGM's internal search/error setting as well as each hybrid's migration controls. \texttt{HybridPGMLIPP} sweeps flush thresholds of \texttt{25000/50000/100000} for the lookup-heavy workload and \texttt{100000/250000/500000} for the insertion-heavy workload. \texttt{HybridPGMLIPPIncremental} uses the same threshold family and additionally sweeps a bounded migration budget per insert, using smaller budgets for the lookup-heavy case and larger budgets for the insertion-heavy case.

\paragraph{Scope.}
The pipeline supports Facebook, Books, and OSMC for both mixed workloads. The required milestone report below presents Facebook, and the second page adds a supplemental Facebook-only comparison between \texttt{HybridPGMLIPP} and \texttt{HybridPGMLIPPIncremental}.

\section*{Results}

\begin{table}[h]
\centering
\scriptsize
\setlength{\tabcolsep}{4pt}
\begin{tabular}{llrrrr}
\toprule
Workload & Index & Throughput (Mops/s) & Size (GiB) & Search & Flush \\
\midrule
Mixed (90% Lookup / 10% Insert) & Dynamic PGM & 1.051 & 1.588 & BinarySearch / 64 & - \\
 & LIPP & 15.648 & 11.829 & - & - \\
 & Hybrid & 1.757 & 11.829 & BinarySearch / 64 & 25000 \\
\addlinespace
Mixed (10% Lookup / 90% Insert) & Dynamic PGM & 3.514 & 1.586 & BinarySearch / 128 & - \\
 & LIPP & 3.246 & 11.787 & - & - \\
 & Hybrid & 1.946 & 11.769 & BinarySearch / 128 & 500000 \\
\bottomrule
\end{tabular}
\caption{Best mean-performing configuration per index for the reported Facebook mixed workloads. Dynamic PGM and the hybrid are selected after averaging three repeats and sweeping their candidate hyperparameters.}
\end{table}


\begin{figure}[h]
\centering
\includegraphics[width=\textwidth]{milestone2_metrics.png}
\caption{Facebook Milestone 2 results. The top row reports throughput and index size for the lookup-heavy mixed workload (90\% lookup, 10\% insertion). The bottom row reports the same metrics for the insertion-heavy mixed workload (10\% lookup, 90\% insertion). Each bar uses the best mean-performing configuration from the three-repeat sweep for that index and workload.}
\end{figure}

\paragraph{Takeaway.}
The required checkpoint comparison uses the best Facebook configuration for Dynamic PGM, LIPP, and the naive hybrid. On this dataset, LIPP remains strongest for the lookup-heavy mixed workload, Dynamic PGM remains strongest for the insertion-heavy mixed workload, and the naive hybrid sits between them while preserving the required combined design.

\clearpage
\section*{Supplemental Hybrid Sweep}

\paragraph{Additional Approach.}
The second approach added beyond the milestone baseline is \texttt{HybridPGMLIPPIncremental}. This section compares the best Facebook configuration of \texttt{HybridPGMLIPP} against the best Facebook configuration of \texttt{HybridPGMLIPPIncremental} under the same mixed workloads.

\begin{table}[h]
\centering
\scriptsize
\setlength{\tabcolsep}{4pt}
\begin{tabular}{p{2.2in}p{1.0in}p{2.2in}}
\toprule
Workload & Hybrid Variant & Best Configuration \\
\midrule
Mixed (90% Lookup / 10% Insert) & \texttt{HybridPGMLIPP}: 2.830 Mops/s & BinarySearch / 128 / flush=25000 \\
 & \texttt{HybridPGMLIPPIncremental}: 1.609 Mops/s & BinarySearch / 64 / flush=25000 / budget=256 \\
\addlinespace
Mixed (10% Lookup / 90% Insert) & \texttt{HybridPGMLIPP}: 2.041 Mops/s & BinarySearch / 512 / flush=500000 \\
 & \texttt{HybridPGMLIPPIncremental}: 1.802 Mops/s & BinarySearch / 512 / flush=500000 / budget=1024 \\
\bottomrule
\end{tabular}
\caption{Best Facebook configurations from the supplemental hybrid-only sweep. \texttt{HybridPGMLIPPIncremental} adds a bounded per-insert migration budget in addition to the flush threshold used by \texttt{HybridPGMLIPP}.}
\end{table}


\begin{figure}[h]
\centering
\includegraphics[width=0.94\textwidth]{milestone2_hybrid_comparison_fb.png}
\caption{Supplemental Facebook comparison between \texttt{HybridPGMLIPP} and \texttt{HybridPGMLIPPIncremental}. For each workload, the plotted bars use the best mean-performing configuration found in the corresponding hyperparameter sweep.}
\end{figure}

\paragraph{Hybrid Observation.}
\texttt{HybridPGMLIPPIncremental} is a more realistic migration policy because it avoids a single large synchronous flush and instead amortizes migration over future insertions. In the Facebook experiments here, however, its best configuration still trails the best \texttt{HybridPGMLIPP} configuration for both mixed workloads. The supplemental sweep is still useful because it shows that migration budget and flush threshold materially affect performance, and it provides a second named hybrid design to build on for the final Task~2 experiments.

\end{document}
